\documentclass[11pt,a4paper]{article}
\usepackage[utf8]{inputenc}
\usepackage[T1]{fontenc}
\usepackage{amsmath,amssymb,amsthm}
\usepackage[margin=2.5cm]{geometry}
\usepackage{hyperref}
\hypersetup{colorlinks=true,linkcolor=blue,citecolor=blue,urlcolor=blue}

\newtheorem{definition}{Definition}
\newtheorem{proposition}{Proposition}
\newtheorem{theorem}{Theorem}
\newtheorem{corollary}{Corollary}
\newtheorem{conjecture}{Conjecture}
\newtheorem{remark}{Remark}

\newcommand{\R}{\mathbb{R}}
\newcommand{\T}{\mathbb{T}}
\newcommand{\UM}{\mathcal{U}_M}
\newcommand{\HH}{\mathcal{H}}

\title{\textbf{Pre-Algebraic Foundations for the CT Framework:\\
Smoothly Undulating Numbers, Complexity Flow,\\
and Algebraic Emergence}\\[6pt]
\large ODS Unified --- Paper P00}

\author{Rub\'en Jord\'a Aracil\\
  \small Independent Researcher\\
  \small \href{https://github.com/rjordaaracil/ods-unified-v2}{github.com/rjordaaracil/ods-unified-v2}\\[4pt]
  \small Series:
  P00 (this paper)
  \,/\, \href{https://doi.org/10.5281/zenodo.18995428}{P01}
  \,/\, \href{https://doi.org/10.5281/zenodo.15078277}{P02}
  \,/\, \href{https://doi.org/10.5281/zenodo.19202918}{P03}
}
\date{March 2026}

\begin{document}
\maketitle

\begin{abstract}
We develop a pre-algebraic foundation for the Clifford--Thomas (CT)
computational framework in which no algebraic structure is assumed
\emph{a priori}. The primitive objects are bounded oscillating series
(smoothly undulating numbers). We introduce a complexity flow
$\partial_\lambda u = G_\lambda[u]$ that preserves global bounds,
admits local complexity increase, and defines limit attractors.
We then define reading modes as integral operators and prove that,
when the continuum of readings collapses to finitely many rigid classes,
the limiting operators satisfy the Clifford anticommutation relations
$\gamma_\mu\gamma_\nu + \gamma_\nu\gamma_\mu = 2g_{\mu\nu}\,\mathrm{Id}$.
Thus Clifford algebra emerges as a boundary state of a purely analytic
theory, rather than being postulated. All definitions are supported by
19 unit tests in the companion code.
\end{abstract}

% ================================================================
\section{Introduction}

The CT runtime~\cite{zenodo_p01,zenodo_p02,zenodo_p03} currently operates
on fields valued in the Clifford algebra $\mathrm{Cl}(3,0)$, an 8-dimensional
graded algebra with generators satisfying
$\gamma_\mu\gamma_\nu + \gamma_\nu\gamma_\mu = 2\delta_{\mu\nu}$.
This algebraic structure is taken as axiomatic in P01--P03.

The present paper asks: \emph{can the Clifford relations be derived
rather than assumed?}

We answer affirmatively by constructing a purely analytic framework
in which:
\begin{enumerate}
\item The only primitive objects are bounded oscillating series
      (no vectors, no matrices, no fixed dimension).
\item A complexity flow generates richness without breaking bounds.
\item Reading modes extract finite observations from the continuum.
\item When readings stabilize at discrete classes, Clifford structure emerges.
\end{enumerate}

No algebra is used until it appears as a consequence.

% ================================================================
\section{Smoothly Undulating Numbers}

Let $I = \T$ be a compact reading parameter. Fix a Schauder basis
$(\phi_n)_{n\ge0} \subset C^1(I)$. Define
\[
\UM := \left\{ u = \sum_{n\ge0} c_n \phi_n \in C^1(I) :
  \|u\|_\infty \le M \right\}.
\]

\begin{definition}[Analytic oscillator]
\label{def:oscillator}
A function $\omega \in C^2(\R)$ is an \emph{analytic oscillator} if:
$\sup|\omega| = 1$, $\inf\omega = -1$,
$\omega(\xi + P) = \omega(\xi)$ for some $P > 0$,
$\int_0^P \omega(\xi)\,d\xi = 0$,
and all critical points of $\omega$ are isolated.
No trigonometry is required.
\end{definition}

\begin{definition}[Smoothly undulating number]
\label{def:sun}
An element $u \in \UM$ is a \emph{smoothly undulating number} (SUN)
if there exists a factorization
\[
u(s) = \alpha(s)\,\omega(\varphi(s)),
\]
with $\alpha \in W^{1,\infty}(I)$, $0 \le \alpha \le M$;
$\varphi \in C^1(I)$, $\varphi'(s) \ge \kappa > 0$;
$\omega$ an analytic oscillator; and the slowness condition
\begin{equation}\label{eq:A1}
|\alpha'(s)| \le \vartheta\,(1 + \alpha(s))\,\varphi'(s),
\qquad 0 \le \vartheta < 1.  \tag{A1}
\end{equation}
The triple $(\alpha, \omega, \varphi)$ is an \emph{admissible factorization},
not the object itself.
\end{definition}

\begin{definition}[Equivalence of factorizations]
\label{def:equiv}
Two factorizations $(\alpha,\omega,\varphi)$ and
$(\tilde\alpha,\tilde\omega,\tilde\varphi)$ represent the same SUN if
there exists $h \in C^1(\R)$, strictly increasing, with $h' > 0$,
such that $\tilde\varphi = h \circ \varphi$,
$\tilde\omega = \omega \circ h^{-1}$, $\tilde\alpha = \alpha$.
\end{definition}

\begin{proposition}[Phase reparametrization invariance]
\label{prop:invariance}
The relation of Definition~\ref{def:equiv} is an equivalence relation
and leaves $u$ invariant.
\end{proposition}
\begin{proof}
$\tilde\alpha(s)\,\tilde\omega(\tilde\varphi(s))
= \alpha(s)\,\omega(h^{-1}(h(\varphi(s))))
= \alpha(s)\,\omega(\varphi(s))
= u(s)$.
Reflexivity, symmetry, and transitivity are immediate. \qed
\end{proof}

\begin{proposition}[Series compatibility]
\label{prop:series}
Every smoothly undulating number belongs to the permitted ontological
class: it is a bounded series.
\end{proposition}
\begin{proof}
By definition $u \in C^1(I) \subset C(I)$. Since $(\phi_n)$ is a
Schauder basis of $C(I)$, there exists a unique expansion
$u = \sum_{n\ge0} c_n \phi_n$ with uniform convergence.
The bound $\|u\|_\infty \le M$ holds by construction. \qed
\end{proof}

\begin{proposition}[Sequences as readings]
\label{prop:sequences}
For any $(b_n)_{n\ge0} \in \ell^\infty$, there exist
$u \in \UM$ with $M = \sup_n|b_n|$ and smooth readers $(L_n)$
such that $L_n[u] = b_n$ for all~$n$.
\end{proposition}
\begin{proof}
Take disjoint intervals $J_n \subset I$ and smooth functions
$\psi_n, \theta_n$ supported in $J_n$ with
$\int_I \psi_n\theta_n = 1$ and $\int_I \psi_m\theta_n = 0$ for $m \ne n$.
Set $u = \sum_n b_n\psi_n$ and $L_n[v] = \int_I v\,\theta_n$.
Then $\|u\|_\infty \le M$ by disjoint support, and $L_n[u] = b_n$. \qed
\end{proof}

% ================================================================
\section{Complexity Flow}

Fix two smoothing operators (heat semigroup):
\[
S_c := e^{\varepsilon_c\partial_{ss}}, \quad
S_f := e^{\varepsilon_f\partial_{ss}}, \quad
0 < \varepsilon_f < \varepsilon_c.
\]
Define the \emph{defect} $A[u] := u - S_c u$,
the \emph{fine return} $\mathcal{R}[u] := S_f u - S_c u$,
the \emph{complexity density} $Q[u] := |A[u]|$,
and the \emph{production proxy} $\Pi[u] := |\mathcal{R}[u]|$.

\begin{definition}[Dynamic permeability]
\label{def:perm}
For bounded continuous $\alpha,\beta,\mu: [0,\infty) \to (0,\infty)$
and $\tau > 0$:
\begin{equation}\label{eq:B1}
p_\lambda[u] := \beta(\lambda)\,
\frac{\Pi[u]}{\Pi[u] + \mu(\lambda)\,Q[u] + \tau}. \tag{B1}
\end{equation}
Then $0 \le p_\lambda[u] \le \beta(\lambda)$, and the local
rupture threshold is $\Pi[u] > \mu(\lambda)\,Q[u] + \tau$.
\end{definition}

\begin{definition}[Complexity flow]
\label{def:flow}
\begin{equation}\label{eq:flow}
\partial_\lambda u = G_\lambda[u] :=
-\alpha(\lambda)\,A[u]
+ p_\lambda[u]\left(1 - \frac{u^2}{M^2}\right)\mathcal{R}[u]. \tag{F}
\end{equation}
\end{definition}

\begin{proposition}[Global existence and boundedness]
\label{prop:bounded}
For any $u_0 \in \UM$, the flow~\eqref{eq:flow} has a unique global
solution in $C^1([0,\infty); C(I))$ satisfying
$\|u(\lambda)\|_\infty \le M$ for all $\lambda \ge 0$.
\end{proposition}
\begin{proof}
$G_\lambda$ is locally Lipschitz in $C(I)$ because the denominator
in~\eqref{eq:B1} is bounded below by $\tau > 0$.
Existence follows from Picard--Lindel\"of in Banach space.

For the bound: if $u(\lambda_*,s_*) = M$ at a first contact point,
then $A[u] = u - S_c u \ge 0$ (since $S_c$ is positive and sub-averaging),
so the first term in~\eqref{eq:flow} is $\le 0$.
The second term vanishes because $1 - u^2/M^2 = 0$.
Thus $\partial_\lambda u \le 0$ at the boundary.
The argument for $-M$ is symmetric.
Therefore $[-M, M]$ is an invariant tube, preventing blow-up. \qed
\end{proof}

\begin{proposition}[Complexity production identity]
\label{prop:production}
Define the complexity energy $\mathcal{E}[u] := \frac{1}{2}\int_I u\,\mathcal{R}[u]\,ds$.
Along the flow~\eqref{eq:flow}:
\begin{equation}\label{eq:B3}
\frac{d}{d\lambda}\mathcal{E}[u(\lambda)]
= -\alpha(\lambda)\int_I A[u]\,\mathcal{R}[u]\,ds
+ \int_I p_\lambda[u]\left(1-\frac{u^2}{M^2}\right)\Pi[u]^2\,ds.
\tag{B3}
\end{equation}
Moreover, $\int_I A[u]\,\mathcal{R}[u]\,ds \ge 0$.
\end{proposition}
\begin{proof}
Since $\mathcal{R}$ is self-adjoint,
$\frac{d}{d\lambda}\mathcal{E} = \int_I \partial_\lambda u\,\mathcal{R}[u]\,ds$.
Substituting~\eqref{eq:flow} and using $|\mathcal{R}[u]| = \Pi[u]$
yields~\eqref{eq:B3}.
For the positivity: $A = I - S_c$ and $\mathcal{R} = S_f - S_c$ are
increasing functions of the same heat operator, hence commute and
are both positive, giving $\int A\,\mathcal{R} \ge 0$. \qed
\end{proof}

\begin{corollary}[Rupture and complexity increase]
\label{cor:rupture}
Define the rupture set
$\Omega_{\mathrm{rup}}(\lambda) := \{s : \Pi[u] > \mu(\lambda)Q[u] + \tau\}$.
On $\Omega_{\mathrm{rup}}$, the permeability satisfies
$p_\lambda[u] > \beta(\lambda)/2$.
If $|u| < M$ on a positive-measure subset of $\Omega_{\mathrm{rup}}$
and local production exceeds dissipation, then
$\frac{d}{d\lambda}\mathcal{E}[u(\lambda)] > 0$.
\end{corollary}

\begin{definition}[Limit attractor]
\label{def:attractor}
Assuming $\alpha(\lambda) \to \alpha_\infty$,
$\beta(\lambda) \to \beta_\infty$,
$\mu(\lambda) \to \mu_\infty$ as $\lambda \to \infty$,
a \emph{limit attractor} is a stable solution $u_* \in \UM$ of
the stationary equation $G_\infty[u_*] = 0$.
\end{definition}

\begin{conjecture}[Generic convergence]
\label{conj:convergence}
For an open set of smoothly undulating initial data and regular
schedules $(\alpha,\beta,\mu)$, the flow converges to a non-empty
family of limit attractors with finite rigid rank.
\end{conjecture}

% ================================================================
\section{Algebraic Emergence}

\begin{definition}[Reading modes]
\label{def:reading}
For a compact parameter set $\Xi$ with $\xi = (\eta,\sigma,\delta)$
and resolution $\varepsilon > 0$, a \emph{reading mode} is an
integral operator
\[
(L_\xi^\varepsilon u)(s) = \int_I K_\xi^\varepsilon(s,r)\,u(r)\,dr,
\]
with smooth kernel $K_\xi^\varepsilon$ satisfying uniform boundedness.
These are \emph{not} ontological layers; they are readers with focus~$\eta$,
scale~$\sigma$, and penumbra~$\delta$.
\end{definition}

\begin{definition}[Quantization radius]
\label{def:quant}
\[
\Delta_N(\varepsilon; u) := \inf_{v_1,\ldots,v_N \in C(I)}
\sup_{\xi \in \Xi} \min_{1\le j\le N}
\|L_\xi^\varepsilon u - v_j\|_\infty.
\]
\end{definition}

\begin{definition}[Rigid rank]
\label{def:rigidrank}
$N_{\mathrm{rig}}(u) := \min\{N \ge 1 :
\lim_{\varepsilon\to0}\Delta_N(\varepsilon;u) = 0\}$,
if it exists. A finite rigid rank means the continuum of readers
collapses to $N$ discrete classes.
\end{definition}

\begin{definition}[Symmetrized closure defect]
\label{def:defect}
For class representatives $L_1^\varepsilon,\ldots,L_N^\varepsilon$
of the $N$ rigid classes, define
\[
D_{\mu\nu}^\varepsilon[u;c] := \sup_{\substack{v \in \HH_\varepsilon[u]\\
\|v\|_\infty \le 1}}
\left\| L_\mu^\varepsilon L_\nu^\varepsilon v
+ L_\nu^\varepsilon L_\mu^\varepsilon v - 2c\,v \right\|_\infty,
\]
where $\HH_\varepsilon[u]$ is the family generated by iterated reading.
Let $g_{\mu\nu}^\varepsilon[u]$ be any minimizer in~$c$, and
$D_{\mu\nu}^\varepsilon[u] := D_{\mu\nu}^\varepsilon[u; g_{\mu\nu}^\varepsilon]$.
\end{definition}

\begin{definition}[Rigidity functional]
\label{def:rigidity}
\[
\mathbf{Rig}_N(\varepsilon; u) :=
\Delta_N(\varepsilon; u) + \max_{1\le\mu,\nu\le N}
D_{\mu\nu}^\varepsilon[u].
\]
First term: discretization. Second term: algebraic closure.
\end{definition}

\begin{definition}[Algebraic emergence of rank $N$]
\label{def:emergence}
$u$ exhibits \emph{algebraic emergence of rank $N$} if
$N_{\mathrm{rig}}(u) = N$ and
$\lim_{\varepsilon\to0}\mathbf{Rig}_N(\varepsilon;u) = 0$.
\end{definition}

\begin{theorem}[Recovery of Clifford relations]
\label{thm:clifford}
Let $u_*$ be a limit attractor of the flow. Suppose:
\begin{enumerate}
\item $N_{\mathrm{rig}}(u_*) = N < \infty$.
\item $\mathbf{Rig}_N(\varepsilon; u_*) \to 0$.
\item For each $\mu$, the strong limit
$\gamma_\mu v := \lim_{\varepsilon\to0} L_\mu^\varepsilon v$
exists for all $v \in \HH[u_*]$.
\item The limit $g_{\mu\nu} := \lim_{\varepsilon\to0}
g_{\mu\nu}^\varepsilon[u_*]$ exists.
\end{enumerate}
Then
\begin{equation}\label{eq:clifford}
\gamma_\mu\gamma_\nu + \gamma_\nu\gamma_\mu
= 2g_{\mu\nu}\,\mathrm{Id}
\quad \text{on } \HH[u_*].
\end{equation}
\end{theorem}
\begin{proof}
By Definition~\ref{def:defect},
$\|L_\mu^\varepsilon L_\nu^\varepsilon v
+ L_\nu^\varepsilon L_\mu^\varepsilon v
- 2g_{\mu\nu}^\varepsilon v\|_\infty \to 0$
for all $v \in \HH_\varepsilon[u_*]$ with $\|v\|_\infty \le 1$.
Passing to the strong limits in the operators and the scalar limit
in $g_{\mu\nu}^\varepsilon$ yields~\eqref{eq:clifford}. \qed
\end{proof}

\begin{corollary}[Squares $\pm 1$]
\label{cor:squares}
If $g_{\mu\mu} \ne 0$, the renormalized reader
$\widetilde\gamma_\mu := |g_{\mu\mu}|^{-1/2}\gamma_\mu$ satisfies
$\widetilde\gamma_\mu^2 = \mathrm{sgn}(g_{\mu\mu})\,\mathrm{Id}$.
If $g_{\mu\nu} = 0$ for $\mu \ne \nu$, then
$\widetilde\gamma_\mu\widetilde\gamma_\nu
+ \widetilde\gamma_\nu\widetilde\gamma_\mu
= 2\eta_{\mu\nu}\,\mathrm{Id}$
with $\eta = \mathrm{diag}(\pm1,\ldots,\pm1)$.
\end{corollary}

\begin{corollary}[$\mathrm{Cl}(3,0)$ sub-closure]
\label{cor:cl30}
If $N_{\mathrm{rig}}(u_*) = 8$ and there exist three classes
$\{\mu_1,\mu_2,\mu_3\}$ with $g_{\mu_i\mu_j} = \delta_{ij}$,
then an emergent sub-closure of type $\mathrm{Cl}(3,0)$ exists
within the attractor's eight reading modes.
This matches the current ODS runtime without assuming eight
coordinates \emph{a priori}.
\end{corollary}

% ================================================================
\section{Experiments}

All definitions and propositions are accompanied by computational
tests in the companion code \texttt{\_claude\_calculo/}.

\begin{table}[ht]
\centering
\begin{tabular}{lcc}
\hline
\textbf{Test group} & \textbf{Tests} & \textbf{Status} \\
\hline
Def 1 (oscillator)     & 3  & 3/3  \\
Def 2 (SUN)            & 2  & 2/2  \\
Prop 1 (invariance)    & 1  & 1/1  \\
Prop 3 (sequences)     & 1  & 1/1  \\
Prop 4 (boundedness)   & 3  & 3/3  \\
Prop 5 (energy)        & 2  & 2/2  \\
Cor 1 (rupture)        & 2  & 2/2  \\
Emergence (Defs 9--14) & 5  & 5/5  \\
\hline
\textbf{Total}         & \textbf{19} & \textbf{19/19} \\
\hline
\end{tabular}
\caption{Test coverage for the pre-algebraic foundations.}
\end{table}

% ================================================================
\section{Discussion}

\subsection{What this paper does}
\begin{itemize}
\item Provides a purely analytic foundation for CT.
\item Derives Clifford structure as a consequence, not an axiom.
\item Defines exact criteria for the continuum-to-discrete transition.
\item Preserves global bounds throughout the complexity flow.
\end{itemize}

\subsection{What this paper does not claim}
\begin{itemize}
\item Generic convergence of the flow to rigid attractors
      (Conjecture~\ref{conj:convergence}, not proven).
\item That the emergent metric $g_{\mu\nu}$ necessarily matches
      any physical metric.
\item That this framework resolves any open problem in PDE theory.
\item That the specific reading mode kernels used in the experiments
      are the ``correct'' or unique choice.
\end{itemize}

% ================================================================
\section{Conclusion}

We have shown that Clifford algebra can emerge from a purely analytic
framework of bounded oscillating series, without postulating any
algebraic structure. The key mechanism is the collapse of a continuum
of reading modes to finitely many rigid classes under resolution
refinement. When this rigidity condition is met (Definition~\ref{def:emergence}),
the limiting operators satisfy the Clifford anticommutation
relations (Theorem~\ref{thm:clifford}).

This provides a foundational justification for the CT runtime
series P01--P06: the $\mathrm{Cl}(3,0)$ structure used there
is not arbitrary but arises as a natural boundary state of the
complexity flow.

The single open problem is Conjecture~\ref{conj:convergence}:
proving generic convergence of the flow to finite-rank rigid
attractors.

\begin{thebibliography}{9}
\bibitem{zenodo_p01}
R.~Jord\'a Aracil,
\emph{CT Runtime: Dual-Clock Clifford Field Integrator},
ODS Unified P01, 2026.
\url{https://doi.org/10.5281/zenodo.18995428}

\bibitem{zenodo_p02}
R.~Jord\'a Aracil,
\emph{Regimes, Soliton Structure, and Inter-Grade Modulation},
ODS Unified P02, 2026.
\url{https://doi.org/10.5281/zenodo.15078277}

\bibitem{zenodo_p03}
R.~Jord\'a Aracil,
\emph{Observability and Stress Monitoring: The $\chi$ Trigger Framework},
ODS Unified P03, 2026.
\url{https://doi.org/10.5281/zenodo.19202918}

\bibitem{ods_code}
R.~Jord\'a Aracil,
\emph{ODS Unified v2 --- source code},
\url{https://github.com/rjordaaracil/ods-unified-v2}
\end{thebibliography}

\end{document}
